\clearpage
\section{Итоговый отчёт}

Заключительный этап учебной практики был посвящён последовательной разработке системы выполнения команд. Сначала был создан отдельный поток с очередью операций и двумя способами завершения. Затем система была дополнена планировщиком, который позволяет выполнять длительные операции небольшими частями. На последнем шаге разработанная схема была проверена на пяти экземплярах команды, каждый из которых должен был получить управление три раза. Связь между заданиями являлась принципиальной: планировщик из задания 18 использовал поток из задания 17, а длительные команды из задания 19 выполнялись средствами этого планировщика.

\begin{table}[h]
    \centering
    \caption{Состав работ заключительного этапа практики}
    \small
    \begin{tabular}{>{\centering\arraybackslash}p{1.1cm} >{\RaggedRight\arraybackslash}p{8.7cm} >{\centering\arraybackslash}p{2cm} >{\centering\arraybackslash}p{2.1cm}}
        \toprule
        \textbf{№} & \textbf{Задание} & \textbf{Тестов} & \textbf{Результат CI} \\
        \midrule
        17 & Разработка многопоточного сервера команд & 10 & успешно \\
        18 & Реализация планировщика команд & 11 & успешно \\
        19 & Реализация длительных команд & 6 & успешно \\
        \bottomrule
    \end{tabular}
\end{table}

\subsection{Разработка многопоточного сервера команд}

Работа над сервером команд началась с определения минимального контракта операции. Интерфейс \codebreak{ICommand} содержит единственный метод \codebreak{Execute}, поэтому поток не зависит от конкретного назначения команды. После этого был создан класс \codebreak{ServerThread}. В его конструкторе формируется отдельный объект \codebreak{Thread}, а метод \codebreak{Start} запускает цикл обработки. Новые команды передаются через потокобезопасную очередь \codebreak{BlockingCollection<ICommand>}, построенную на стандартных средствах .NET \cite{threadsafe}. Такой выбор позволил отделить добавление команд внешними потоками от их последовательного исполнения серверным потоком.

Следующим шагом стало определение правил остановки. Команда \codebreak{HardStopCommand} переводит сервер в режим немедленного завершения. После её выполнения цикл прекращается, даже если в очереди остаются элементы. Команда \codebreak{SoftStopCommand}, напротив, запрещает дальнейшее добавление операций, но сохраняет обработку уже принятых элементов. Поток заканчивает работу только после опустошения очереди. Для исключения ошибочного управления предусмотрена проверка текущего потока: обе команды остановки могут успешно выполниться только внутри того экземпляра \codebreak{ServerThread}, для которого они созданы. Попытка вызвать их извне приводит к исключению.

Основное затруднение заключалось в различии между немедленным и корректным отложенным завершением. Простого логического признака остановки было недостаточно. Если устанавливать его одинаково для обоих режимов, то мягкая остановка теряет оставшиеся команды, а жёсткая, наоборот, продолжает ненужную обработку. Проблема была решена введением двух состояний завершения и отдельной проверкой режима после каждого вызова \codebreak{Execute}. Для мягкой остановки используется завершение добавления элементов в \codebreak{BlockingCollection}; перечисление существующих элементов при этом продолжается до конца.

Второе затруднение было связано с поведением пустой очереди. Цикл с постоянной проверкой \codebreak{TryTake} работал бы даже при полном отсутствии полезной нагрузки и расходовал процессорное время пропорционально длительности простоя. В реализации применено блокирующее извлечение. Если элементов нет, системный поток переводится в состояние ожидания и возобновляется после добавления команды. Такой способ соответствует назначению серверного потока: ресурсы используются только при появлении работы, а задержка ожидания не создаёт активной нагрузки на процессор \cite{threads}.

Исключение внутри одной команды не должно разрушать весь цикл обработки. Поэтому вызов \codebreak{Execute} помещён в блок обработки исключений. При наличии \codebreak{IExceptionHandler} в него передаются сама команда и полученное исключение, после чего сервер переходит к следующему элементу. Отдельными тестами проверялись оба вида остановки, запрет их внешнего вызова, обработка исключения, ожидание новой команды при пустой очереди и одновременная подача операций несколькими производителями. Наиболее сложным моментом задания оказалось согласование состояния очереди и режима завершения так, чтобы жёсткая остановка была немедленной, мягкая не теряла принятые операции, а пустой сервер не выполнял бесполезный цикл.

\subsection{Реализация планировщика команд}

Во втором задании потребовалось устранить монопольное использование потока длительной операцией. Если обычная команда выполняет всю работу за один вызов, сервер не может передать управление другим элементам до возврата из \codebreak{Execute}. Для изменения этого поведения длительная операция была разделена на небольшие шаги. Интерфейс \codebreak{ILongRunningCommand} расширяет \codebreak{ICommand} свойством \codebreak{IsCompleted}. Один вызов выполняет только очередной шаг, а признак завершения сообщает серверу, требуется ли планировать следующий вызов.

Планировщик определён интерфейсом \codebreak{IScheduler}, который содержит операции проверки наличия работы, выбора следующей команды и добавления команды. Реализация \codebreak{RoundRobinScheduler} хранит элементы в очереди. Выбранная команда удаляется из начала, выполняет один шаг и при незавершённом состоянии добавляется в конец. Благодаря этому при наличии команд $A$, $B$ и $C$ образуется последовательность $A,B,C,A,B,C,\ldots$. Каждая операция регулярно получает управление, поэтому одна длительная задача не вытесняет остальные.

Главная трудность возникла при совместном ожидании двух источников работы. Новые команды продолжают поступать во входную очередь, а ранее начатые длительные операции находятся в планировщике. Блокирующий вызов \codebreak{Take} для входной очереди использовать нельзя: поток может заснуть при пустой входной очереди, хотя планировщик уже содержит незавершённую работу. Постоянный неблокирующий опрос обеих структур также неприемлем, поскольку при полном отсутствии команд он превращается в активный цикл.

Решение основано на сочетании \codebreak{ConcurrentQueue<ICommand>} и \codebreak{SemaphoreSlim} \cite{semaphore}. Из входной очереди элементы извлекаются неблокирующим способом. Если они отсутствуют, поток проверяет планировщик. Ожидание на семафоре начинается только тогда, когда оба источника пусты. Добавление новой команды увеличивает счётчик семафора и пробуждает обработчик. Повторное выполнение длительной операции производится через планировщик, а не через входную очередь, поэтому серверный поток не записывает в ту же очередь, из которой читает. Это устраняет описанный в задании риск мёртвой блокировки при ограниченной вместимости очереди.

Для сохранения отзывчивости поток чередует обработку нового элемента и шага запланированной операции. Несколько длительных команд, поступивших подряд, сначала регистрируются в планировщике, после чего выполняются циклически. Если во время их работы появляется обычная команда, она получает управление между двумя шагами длительной операции. Исключение в одном шаге передаётся обработчику и не приводит к бесконечному повторению неисправной команды. Команды \codebreak{HardStopCommand} и \codebreak{SoftStopCommand} сохраняют прежнюю семантику: жёсткая остановка прекращает цикл, а мягкая ожидает завершения входной очереди и всех элементов планировщика.

Набор тестов проверял порядок Round Robin, выбор из пустого планировщика, выполнение обычной команды до завершения длительной, поступление новой команды между двумя шагами, оба способа остановки, ожидание при полном отсутствии работы и обработку исключения. Наиболее сложной частью оказалось не само циклическое хранение элементов, а организация сна потока. Необходимо было исключить одновременно две противоположные ошибки: засыпание при наличии работы в планировщике и активный опрос при отсутствии любой работы.

\subsection{Реализация длительных команд}

В заключительном задании создан класс \codebreak{TestCommand}. Он хранит идентификатор, счётчик вызовов и требуемое число шагов. После каждого вызова \codebreak{Execute} счётчик увеличивается, а в поток вывода записывается строка вида «Поток 1 вызов 1». При достижении трёх вызовов свойство \codebreak{IsCompleted} получает значение \codebreak{true}. Сервер больше не возвращает такой экземпляр в планировщик. Дополнительный вызов завершённой команды не изменяет счётчик и не создаёт повторной строки.

Для демонстрации сформировано пять экземпляров с идентификаторами от 1 до 5. Все они добавляются в \codebreak{ServerThread} и регистрируются планировщиком. Получается три полных цикла: в первом каждая команда вызывается один раз, во втором счётчики увеличиваются до двух, в третьем до трёх. Консольный проект \codebreak{task19app} выводит ровно 15 строк. Тем самым длительные операции завершаются, но ни одна из них не удерживает серверный поток на протяжении всех трёх шагов.

Отдельная трудность была связана с требованием завершить демонстрацию именно командой \codebreak{HardStop}. Если поместить обычную команду жёсткой остановки во входную очередь сразу после пяти длительных операций, сервер зарегистрирует длительные элементы в планировщике и выполнит остановку раньше их первого шага. Ожидание завершения в главном потоке с постоянной проверкой счётчиков, в свою очередь, повторило бы проблему активного опроса.

Для решения создана длительная команда \codebreak{HardStopAfterCommandsCommand}. Она также участвует в Round Robin и после каждого цикла проверяет признаки завершения пяти рабочих команд. Пока хотя бы один элемент не завершён, команда остаётся в планировщике. После третьего цикла все признаки получают значение \codebreak{true}; тогда команда выполняет \codebreak{HardStop} непосредственно в серверном потоке. Проверка принадлежности потоку проходит успешно, и цикл прекращается. Такой вариант не требует ожидания со стороны консольного приложения и использует тот же механизм планирования, что и основные операции.

Тестами подтверждены три вызова одного экземпляра, отсутствие четвёртого вызова после завершения, поддержка другого числа шагов, проверка ошибочных параметров и полная демонстрация пяти команд. Интеграционный тест сравнивает весь вывод с ожидаемой последовательностью из 15 строк и убеждается, что сервер завершён жёсткой остановкой. Наиболее сложным оказался выбор момента \codebreak{HardStop}: остановка должна выполняться не после фиксированной задержки, а после фактического завершения всех операций и только в целевом потоке.

\subsection{Описание и интерпретация результатов экспериментов}

Эксперименты выполнялись на той же реализации, которая использовалась в заданиях 17--19. Исходные результаты сохранены в файлах формата CSV. Первый опыт оценивал затраты процессорного времени при отсутствии команд. Для каждого варианта проводилось пять запусков длительностью одна секунда. В варианте с блокирующим ожиданием запускался пустой \codebreak{ServerThread}; после одной секунды в него передавалась команда \codebreak{HardStop}. Для сравнения выполнялся активный цикл \codebreak{SpinWait} той же продолжительности. Измерялось изменение \codebreak{Process.TotalProcessorTime}. Перед серией выполнялся короткий прогрев, уменьшающий влияние первичной компиляции методов.

\begin{table}[ht]
    \centering
    \caption{Процессорное время за одну секунду отсутствия полезной работы}
    \small
    \begin{tabular}{>{\centering\arraybackslash}p{2.2cm} >{\centering\arraybackslash}p{5cm} >{\centering\arraybackslash}p{5cm}}
        \toprule
        \textbf{Запуск} & \textbf{Блокирующее ожидание, мс} & \textbf{Активный опрос, мс} \\
        \midrule
        1 & $0{,}000$ & $984{,}375$ \\
        2 & $0{,}000$ & $1000{,}000$ \\
        3 & $0{,}000$ & $984{,}375$ \\
        4 & $0{,}000$ & $1000{,}000$ \\
        5 & $46{,}875$ & $1000{,}000$ \\
        \midrule
        Среднее & $9{,}375$ & $993{,}750$ \\
        \bottomrule
    \end{tabular}
\end{table}

\begin{figure}[ht]
    \centering
    \includegraphics[width=0.94\textwidth]{week4-idle-cpu.png}
    \caption{Сравнение затрат процессорного времени при отсутствии команд}
    \label{fig:week4-idle}
\end{figure}

На рисунке~\ref{fig:week4-idle} столбцы показывают средние значения, а отрезки погрешности соответствуют выборочному стандартному отклонению. Для блокирующего ожидания среднее процессорное время составило $9{,}375$ мс, для активного опроса --- $993{,}750$ мс. Относительное уменьшение равно $99{,}06\%$. Нулевые значения в четырёх запусках означают, что затраты оказались ниже разрешения системного счётчика. Значение $46{,}875$ мс в пятом запуске связано с кратковременной фоновой работой процесса и увеличивает стандартное отклонение, но не меняет общий вывод.

Полученная разница объясняется состояниями системного потока. При вызове \codebreak{SemaphoreSlim.Wait} поток не получает кванты процессорного времени до появления сигнала. При активном опросе условие проверяется непрерывно, поэтому одно логическое ядро остаётся занятым почти всю секунду. Результат подтверждает выполнение требования об отсутствии значительных затрат, зависящих от продолжительности простоя. Небольшие постоянные расходы на запуск, пробуждение и остановку потока допустимы, поскольку они не растут пропорционально времени ожидания.

Во втором опыте фиксировался фактический порядок вызовов пяти экземпляров \codebreak{TestCommand}. Данные получены из консольного вывода без перестановки строк. По оси абсцисс на рисунке~\ref{fig:week4-round-robin} указан общий номер вызова \codebreak{Execute}, по оси ординат --- идентификатор выбранной команды.

\begin{figure}[ht]
    \centering
    \includegraphics[width=0.94\textwidth]{week4-round-robin.png}
    \caption{Порядок выполнения пяти длительных команд планировщиком Round Robin}
    \label{fig:week4-round-robin}
\end{figure}

График состоит из трёх одинаковых циклов $1,2,3,4,5$. После команды с идентификатором 5 планировщик возвращается к идентификатору 1. Ни одна команда не вызывается два раза подряд, пока другие готовые элементы ожидают выполнения. Расстояние между соседними вызовами одного экземпляра постоянно и равно пяти диспетчеризациям. Следовательно, стратегия обеспечивает слабую справедливость: при сохранении команды в планировщике она обязательно получает следующий шаг после конечного числа выборов.

Третий цикл заканчивается вызовом команды 5 с общим номером 15. К этому моменту счётчики всех пяти экземпляров равны трём. Следующая проверка \codebreak{HardStopAfterCommandsCommand} обнаруживает завершение набора и останавливает сервер. На графике отсутствуют четвёртый цикл и преждевременное окончание, что одновременно подтверждает корректность признака \codebreak{IsCompleted}, повторного добавления в конец очереди и условия жёсткой остановки. Форма графика определяется не продолжительностью отдельных шагов, а выбранной дисциплиной обслуживания: очередь Round Robin сохраняет порядок поступления и после каждого незавершённого шага переносит команду в конец.
